\section{Finite separable extensions}

Fix an adelic structure $S=(\Omega,\mathcal A,\nu,\phi)$ on a field $K$.
For an extension $L/K$, use the parameter space $\Omega_L$, projection
$\pi$, evaluations $e_b(x)=|b|_x$, and fibres $X_\omega$ of
Definition~\ref{def:extension-parameters}. Except where explicitly stated
otherwise, $L/K$ is finite separable and $n=[L:K]$. This section follows
\cite[\S3.3]{chen-moriwaki}; the bounded-generator argument and the passage
from measurable to almost everywhere measurable integrands are made explicit.

\subsection{Finite fibres and their integrals}

\begin{lemma}[Finiteness of extension fibres]\label{lem:finite-extension-fibres}
\leanok
\source{chen-moriwaki-adelic:page-0194,chen-moriwaki-adelic:page-0204}
\uses{def:extension-parameters,lem:finite-values-via-embeddings,lem:extension-coordinates}
\lean{ArakelovGeometry.AdelicStructure.ParameterSpaceExtension.finite_extensionFiber}
For every finite extension $L/K$, separable or not, and every $\omega\in\Omega$,
the set $X_\omega$ is finite.
\end{lemma}
\begin{proof}
Sending $(\omega,w)$ to $w$ identifies the fibre with the set of absolute
values on $L$ extending $\phi_\omega$. Lemma~\ref{lem:finite-values-via-embeddings}
realizes these as pullbacks along the finitely many embeddings into a fixed
algebraic closure of $K_{\phi_\omega}$. It applies equally to inseparable
finite extensions.
\end{proof}

\begin{lemma}[Normalization of local degrees]\label{lem:fibre-normalization}
\leanok
\source{chen-moriwaki-adelic:page-0194}
\uses{def:weighted-fibre-integral,thm:local-degree-sum,lem:finite-extension-fibres}
\lean{ArakelovGeometry.AdelicStructure.ParameterSpaceExtension.tsum_fibreWeight}
The weights $p_\omega(x)=[L_x:K_\omega]/n$ satisfy
$\sum_{x\in X_\omega}p_\omega(x)=1$.
\end{lemma}
\begin{proof}
Theorem~\ref{thm:local-degree-sum} gives
$n=\sum_x[L_x:K_\omega]$. Division by the positive integer $n$
gives the assertion.
\end{proof}

\begin{lemma}[Probability and integrability on finite fibres]
\label{lem:finite-fibre-integrability}
\leanok
\lean{ArakelovGeometry.AdelicStructure.ParameterSpaceExtension.isProbabilityMeasure_fibreProbabilityMeasure,ArakelovGeometry.AdelicStructure.ParameterSpaceExtension.integrable_fibre}
\source{chen-moriwaki-adelic:page-0194,chen-moriwaki-adelic:page-0195}
\uses{def:weighted-fibre-integral,lem:finite-extension-fibres,lem:fibre-separation,lem:fibre-normalization}
The measure $P_\omega$ is a probability measure, and every real-valued
function on $X_\omega$ is integrable for $P_\omega$.
\end{lemma}
\begin{proof}
The mass of $P_\omega=\sum_x p_\omega(x)\delta_x$ is
$\sum_x p_\omega(x)=1$. The finite fibre has a separating measurable
structure: intersecting, for each fixed point, measurable sets which
separate it from the other finitely many points shows that its singleton
is measurable. Hence all functions on the fibre are measurable. Such
a real-valued function is bounded, and the integral of its absolute
value is bounded by its maximum absolute value times $P_\omega(X_\omega)=1$.
\end{proof}

\begin{lemma}[Finite weighted sums]\label{lem:finite-fibre-sums}
\leanok
\source{chen-moriwaki-adelic:page-0194,chen-moriwaki-adelic:page-0195}
\uses{def:weighted-fibre-integral,lem:finite-extension-fibres,lem:fibre-separation,lem:finite-fibre-integrability}
\lean{ArakelovGeometry.AdelicStructure.ParameterSpaceExtension.fibreIntegral_eq_sum,ArakelovGeometry.AdelicStructure.ParameterSpaceExtension.fibreLIntegral_eq_sum}
For every real function $f$ and extended nonnegative function $g$ on
$\Omega_L$,
\[
 If(\omega)=\sum_{x\in X_\omega}p_\omega(x)f(x),\qquad
 I^+g(\omega)=\sum_{x\in X_\omega}p_\omega(x)g(x).
\]
\end{lemma}
\begin{proof}
The fibre is finite and its measurable structure separates points, so
each singleton is measurable. The measure is
$P_\omega=\sum_{x\in X_\omega}p_\omega(x)\delta_x$.
Every real function on the fibre is integrable by
Lemma~\ref{lem:finite-fibre-integrability}.
Integrating against the finite sum of measures and using
$\int h\,d\delta_x=h(x)$ gives the real formula. Additivity of the
nonnegative integral with respect to the measure gives the second
formula, including when one of its terms is infinite.
\end{proof}

\begin{lemma}[Limits of finite weighted sums]\label{lem:finite-fibre-limits}
\leanok
\source{chen-moriwaki-adelic:page-0195}
\uses{lem:finite-fibre-sums}
\lean{ArakelovGeometry.AdelicStructure.ParameterSpaceExtension.tendsto_fibreIntegral,ArakelovGeometry.AdelicStructure.ParameterSpaceExtension.fibreLIntegral_iSup,ArakelovGeometry.AdelicStructure.ParameterSpaceExtension.fibreLIntegral_tsum}
The operator $I$ preserves pointwise limits of real functions, and
$I^+$ preserves increasing limits and countable sums of nonnegative
functions.
\end{lemma}
\begin{proof}
For a fixed $\omega$, all the sums in
Lemma~\ref{lem:finite-fibre-sums} are finite.
Limits of real sequences commute with a fixed finite linear
combination. For nonnegative extended functions, multiplication by
a finite nonnegative constant and a finite sum commute with increasing
limits. The same finite sum commutes with countable sums by
reordering a family of nonnegative terms. These facts give each
assertion at $\omega$, and $\omega$ is arbitrary.
\end{proof}

\begin{lemma}[Integral of one]\label{lem:fibre-one}
\leanok
\source{chen-moriwaki-adelic:page-0194,chen-moriwaki-adelic:page-0195}
\uses{lem:finite-fibre-integrability}
\lean{ArakelovGeometry.AdelicStructure.ParameterSpaceExtension.fibreIntegral_one}
$I1=1$.
\end{lemma}
\begin{proof}
The integral of the constant function $1$ is the total mass
of the probability measure $P_\omega$, which is $1$.
\end{proof}

\begin{lemma}[Additivity]\label{lem:fibre-additivity}
\leanok
\source{chen-moriwaki-adelic:page-0195}
\uses{lem:finite-fibre-integrability}
\lean{ArakelovGeometry.AdelicStructure.ParameterSpaceExtension.fibreIntegral_add}
For all real functions $f,g$ on $\Omega_L$, $I(f+g)=If+Ig$.
\end{lemma}
\begin{proof}
Both restrictions to $X_\omega$ are integrable by
Lemma~\ref{lem:finite-fibre-integrability}. Additivity of the real integral
therefore gives the assertion at every $\omega$.
\end{proof}

\begin{corollary}[Retraction of base pullback]\label{cor:fibre-pullback}
\leanok
\source{chen-moriwaki-adelic:page-0196}
\uses{lem:finite-fibre-integrability}
\lean{ArakelovGeometry.AdelicStructure.ParameterSpaceExtension.fibreIntegral_comp_proj}
For every real function $\varphi$ on $\Omega$, $I(\varphi\circ\pi)=\varphi$.
\end{corollary}
\begin{proof}
On $X_\omega$, the function $\varphi\circ\pi$ is the constant
$\varphi(\omega)$. Its integral against the probability measure
$P_\omega$ is $\varphi(\omega)$.
\end{proof}

\begin{proposition}[Compatibility with intermediate fields]\label{prop:fibre-restriction}
\leanok
\source{chen-moriwaki-adelic:page-0195,chen-moriwaki-adelic:page-0196}
\uses{lem:finite-fibre-sums,lem:weighted-restriction-sum,lem:restriction-transitive}
\lean{ArakelovGeometry.AdelicStructure.ParameterSpaceExtension.fibreIntegral_comp_restrict}
If $K\subseteq F\subseteq L$, then for every real function $f$ on
$\Omega_F$,
\[
 I_{L/K}(f\circ\pi_{L/F})=I_{F/K}f.
\]
\end{proposition}
\begin{proof}
Fix $\omega$. Identify the fibres over $\omega$ with the sets of
absolute values extending $\phi_\omega$. Lemma~\ref{lem:weighted-restriction-sum}
applied to $f$ on the intermediate fibre gives
\[
 \sum_{y\in X_{L,\omega}}[L_y:K_\omega]f(y|_F)
 =[L:F]\sum_{x\in X_{F,\omega}}[F_x:K_\omega]f(x).
\]
Divide by $[L:K]=[F:K][L:F]$, using positivity of $[L:F]$ to cancel
that factor. By Lemma~\ref{lem:finite-fibre-sums}, the resulting
identity is the asserted equality of fibre integrals at $\omega$.
\end{proof}

\subsection{Power sums and measurable conjugate averages}

\begin{lemma}[Recovering powers from consecutive power sums]
\label{lem:vandermonde-power-recovery}
\leanok
\lean{ArakelovGeometry.exists_pow_eq_sum_shifted_sum_pow}
\source{chen-moriwaki-adelic:page-0196,chen-moriwaki-adelic:page-0197}
Let $\alpha_1,\ldots,\alpha_d$ be pairwise distinct elements of a field,
and put $s_N=\sum_i\alpha_i^N$. There are coefficients $c_{ir}$,
independent of $N$, such that
\[
 \alpha_i^N=\sum_{r=0}^{d-1}c_{ir}s_{N+r}
 \quad\text{for every }i\text{ and every }N\geq0.
\]
\end{lemma}
\begin{proof}
Let $V$ have entries $V_{ri}=\alpha_i^r$ for $0\leq r<d$.
The Vandermonde determinant is
$\prod_{i<j}(\alpha_j-\alpha_i)$, which is nonzero by distinctness.
Set $c_{ir}=(V^{-1})_{ir}$. If $u_i=\alpha_i^N$, then
$(Vu)_r=\sum_i\alpha_i^{N+r}=s_{N+r}$. Multiplying by $V^{-1}$
gives the asserted identities.
\end{proof}

\begin{lemma}[Decay of powers from decay of power sums]
\label{lem:power-sum-decay}
\leanok
\lean{ArakelovGeometry.tendsto_pow_zero_of_tendsto_sum_pow_zero}
\uses{lem:vandermonde-power-recovery}
\source{chen-moriwaki-adelic:page-0196,chen-moriwaki-adelic:page-0197}
Let $\alpha_1,\ldots,\alpha_d$ be pairwise distinct elements of a
normed field. If $\sum_i\alpha_i^N\to0$ as $N\to\infty$, then
$\alpha_i^N\to0$ for each $i$.
\end{lemma}
\begin{proof}
Each shifted sequence $s_{N+r}$ tends to zero. The fixed finite linear
combination in Lemma~\ref{lem:vandermonde-power-recovery} therefore
tends to zero, proving the claim.
\end{proof}

\begin{lemma}[Recovering the largest absolute value]\label{lem:power-sum-radius}
\leanok
\source{chen-moriwaki-adelic:page-0196,chen-moriwaki-adelic:page-0197}
\uses{lem:power-sum-decay}
\lean{AbsoluteValue.iSup_eq_limsup_sum_pow_rpow}
For a real absolute value $v$ on a field and pairwise distinct elements
$\alpha_1,\ldots,\alpha_d$, where $d\geq1$, one has
\[
 \max_i v(\alpha_i)
 =\limsup_{N\to\infty}v\left(\sum_{i=1}^d\alpha_i^N\right)^{1/N}.
\]
\end{lemma}
\begin{proof}
Put $M=\max_i v(\alpha_i)$ and $s_N=\sum_i\alpha_i^N$.
For $N\geq1$, the triangle inequality gives $v(s_N)\leq dM^N$,
so $v(s_N)^{1/N}\leq d^{1/N}M$. Since $d^{1/N}\to1$, this
proves the upper bound and shows that the root sequence is bounded.
Its terms are nonnegative. If $M=0$, the upper bound proves equality.

Suppose $M>0$. If the limsup were less than $M$, choose a real number
$b$ strictly between that limsup and $M$, with $b>0$. Eventually
$v(s_N)^{1/N}\leq b$, hence $v(s_N)\leq b^N$.
Choose $j$ with $v(\alpha_j)=M$ and put $\beta_i=\alpha_i/\alpha_j$.
The elements $\beta_i$ are pairwise distinct, $\beta_j=1$, and
\[
 v\left(\sum_i\beta_i^N\right)
 =v(s_N)/M^N\leq (b/M)^N
\]
for all sufficiently large $N$. Since $0<b/M<1$, the right side
tends to zero. Equip the field with the norm defined by $v$ and apply
Lemma~\ref{lem:power-sum-decay}. It implies $\beta_j^N\to0$,
contradicting $\beta_j^N=1$. This proves the lower bound.
The argument uses the full Vandermonde matrix over the original field.
\end{proof}

\begin{lemma}[A separable orbit as an embedding set]
\label{lem:separable-orbit-embeddings}
\leanok
\lean{ArakelovGeometry.AdelicStructure.adjoin_restrict_aut_surjective,ArakelovGeometry.AdelicStructure.adjoin_embeddings_gen_injective}
\uses{lem:algebraic-embedding-extension}
\source{chen-moriwaki-adelic:page-0198}
Let $c$ be separable over $K$ in an algebraic closure $C$, and let
$E=K(c)$. Restriction from $\operatorname{Aut}_K(C)$ onto the set of
$K$-embeddings $E\hookrightarrow C$ is surjective, and evaluation at
$c$ is injective on that embedding set.
\end{lemma}
\begin{proof}
An embedding is determined by its value at the generator $c$, which
proves injectivity of evaluation. Every embedding $E\to C$ extends
to an embedding $C\to C$ by Lemma~\ref{lem:algebraic-embedding-extension}.
Its image is algebraically closed, and $C$ is algebraic over that image
because the image contains $K$. An algebraically closed field has no
proper algebraic extension, so this extended embedding is surjective.
It is therefore an automorphism restricting to the given embedding.
\end{proof}

\begin{lemma}[Counting a separable orbit]\label{lem:separable-orbit-cardinality}
\notready
\uses{lem:separable-orbit-embeddings}
\source{chen-moriwaki-adelic:page-0198}
For a separable element $c$ over $K$ in an algebraic closure $C$, the
orbit $\{\tau(c):\tau\in\operatorname{Aut}_K(C)\}$ is the set of roots
of the minimal polynomial $P_c$ and has $[K(c):K]$ elements.
\end{lemma}
\begin{proof}
Every image of $c$ under a $K$-embedding is a root of $P_c$.
Conversely, each root defines an embedding by evaluation on
$K[T]/(P_c)\simeq K(c)$. Lemma~\ref{lem:separable-orbit-embeddings}
identifies these embeddings with distinct orbit elements. The polynomial
$P_c$ splits in $C$ and has distinct roots by separability, so their
number is $\deg P_c=[K(c):K]$.
\end{proof}

\begin{lemma}[Power sums of a separable orbit lie in the base]\label{lem:orbit-power-sums}
\leanok
\source{chen-moriwaki-adelic:page-0198}
\uses{lem:separable-orbit-embeddings}
\lean{ArakelovGeometry.AdelicStructure.sum_pow_embeddings_eq_trace}
Let $E/K$ be a finite separable extension, $C$ an algebraic closure
of $K$, and $b\in E$. For every integer $N\geq0$,
\[
 \sum_{\sigma:E\hookrightarrow C}\sigma(b)^N
 =\operatorname{Tr}_{E/K}(b^N)\in K\subseteq C.
\]
In particular, if $E=K(c)$ and $\alpha_1,\ldots,\alpha_d$ are the
distinct conjugates of $c$, then
$\sum_i\alpha_i^N=\operatorname{Tr}_{K(c)/K}(c^N)$.
\end{lemma}
\begin{proof}
Choose a primitive element $\theta$ of the finite separable extension
$E/K$. Its separable minimal polynomial factors in $C[T]$ into
distinct linear factors.
The Chinese remainder theorem gives
$C\otimes_K E\simeq\prod_{\sigma:E\hookrightarrow C}C$.
Multiplication by $b^N$ on the right is diagonal, with entries
$\sigma(b)^N$, so its trace is their sum. The trace of a matrix is
unchanged by extending scalars, since it is the sum of its diagonal
entries. Thus this sum is the image of $\operatorname{Tr}_{E/K}(b^N)$.
For $E=K(c)$ and $b=c$, Lemma~\ref{lem:separable-orbit-embeddings}
identifies the summands with the distinct conjugates. This trace identity
is the power-sum case of the symmetric-polynomial argument in the
reference. It involves no division by an integer and applies in
positive characteristic as well.
\end{proof}

\begin{lemma}[Measurability of orbit maxima]\label{lem:orbit-maximum}
\leanok
\source{chen-moriwaki-adelic:page-0196,chen-moriwaki-adelic:page-0198}
\uses{def:adelic-structure,lem:power-sum-radius,lem:orbit-power-sums,lem:separable-orbit-embeddings,lem:separable-orbit-cardinality,def:extension-parameters,lem:extension-coordinates}
\lean{ArakelovGeometry.AdelicStructure.measurable_iSup_abv_aut}
Choose for each $\omega$ any extension $v_\omega$ of $\phi_\omega$ to
an algebraic closure $\overline K$. For every separable $c\in\overline K$,
\[
 \omega\longmapsto
 \sup_{\tau\in\operatorname{Aut}_K(\overline K)}v_\omega(\tau c)
\]
is measurable; no measurability of the choices $v_\omega$ is required.
\end{lemma}
\begin{proof}
The orbit of $c$ is the finite set of distinct roots of its minimal
polynomial: each embedding of $K(c)$ extends to an automorphism of the
algebraic closure. For every $N\geq1$, its power sum $a_N$ lies in $K$
by Lemma~\ref{lem:orbit-power-sums}. If $a_N=0$, then
$v_\omega(a_N)=0$ is measurable. Otherwise positivity gives
\[
 v_\omega(a_N)=\phi_\omega(a_N)
 =\exp\bigl(\log\phi_\omega(a_N)\bigr),
\]
which is measurable by the adelic hypothesis. Composing with the continuous
map $t\mapsto t^{1/N}$ on $[0,\infty)$ preserves measurability.
Lemma~\ref{lem:power-sum-radius} identifies the desired maximum with
the limsup of this sequence. The limsup is a countable infimum of
countable suprema of measurable extended real functions, and its values
are finite, so it is a measurable real function.
\end{proof}

\begin{lemma}[Counting embeddings above a place]\label{lem:conjugate-count}
\source{chen-moriwaki-adelic:page-0198,chen-moriwaki-adelic:page-0199}
\uses{lem:local-factorization,lem:complete-extension}
\notready
Fix $\omega$ and realize $\overline K$ inside an algebraic closure of
$K_\omega$, with its extended absolute value. Among the $n$ embeddings
$\sigma:L\hookrightarrow\overline K$ over $K$, the number inducing
a given $x\in X_\omega$ is $[L_x:K_\omega]$. Consequently, for $b\in L$,
\[
 I_{L/K}e_b(\omega)=\frac1n\sum_{\sigma:L\hookrightarrow\overline K}
 |\sigma b|_\omega.
\]
\end{lemma}
\begin{proof}
Every $K$-embedding of $L$ into the algebraic closure of $K_\omega$
extends uniquely to a $K_\omega$-algebra homomorphism on
$K_\omega\otimes_KL$, by $a\otimes b\mapsto a\sigma(b)$.
This algebra is the product of the separable fields $L_x$.
A homomorphism from this product to an algebraically closed field
selects exactly one factor: its idempotents have images zero or one,
and their sum is one. A separable factor $L_x$ has exactly
$[L_x:K_\omega]$ embeddings into the algebraic closure of $K_\omega$.
Each gives the absolute value $x$ on $L$, by uniqueness of extension
of the absolute value from the complete field $K_\omega$ to its algebraic
extensions. Distinct factors correspond to distinct absolute values,
so the embeddings inducing $x$ are precisely those selecting $L_x$.
The subfield of elements algebraic over $K$ in the chosen algebraic
closure of $K_\omega$ is exactly the image of $\overline K$: an element
algebraic over $K$ is algebraic over that image, and an algebraically
closed field has no proper algebraic field extension. Thus these are
exactly the embeddings $L\hookrightarrow\overline K$ in the statement.
For every such embedding selecting $L_x$, the absolute value of
$\sigma b$ is $|b|_x$. Grouping the $n$ embeddings by the selected
factor therefore gives
\[
 \frac1n\sum_\sigma|\sigma b|_\omega
 =\sum_{x\in X_\omega}\frac{[L_x:K_\omega]}n|b|_x
 =I_{L/K}e_b(\omega).
\]
\end{proof}

\begin{lemma}[Largest products in an ordered tuple]
\label{lem:ordered-subset-product}
\leanok
\lean{ArakelovGeometry.prod_le_prod_initial_of_antitone,ArakelovGeometry.prod_initial_eq_iSup_prod,ArakelovGeometry.prod_initial_comp_perm_eq_iSup_prod}
\source{chen-moriwaki-adelic:page-0198}
Let $s_0\geq\cdots\geq s_{n-1}\geq0$, and let $0\leq k\leq n$.
For every $k$-element subset $J\subseteq\{0,\ldots,n-1\}$,
\[
 \prod_{j\in J}s_j\leq\prod_{i=0}^{k-1}s_i.
\]
Consequently the right side is the largest product over all such subsets.
\end{lemma}
\begin{proof}
Enumerate $J$ increasingly as $j_0<\cdots<j_{k-1}$. Induction gives
$j_i\geq i$: the base case is nonnegativity of $j_0$, and
$j_{i+1}\geq j_i+1\geq i+1$. Since the tuple is decreasing,
$s_{j_i}\leq s_i$. Multiplying these inequalities between nonnegative
numbers proves the bound. The subset $\{0,\ldots,k-1\}$ attains it.
Permuting an arbitrary nonnegative tuple into decreasing order preserves
the collection of its $k$-subset products, so the same maximum describes
the product of its $k$ largest entries.
\end{proof}

\begin{lemma}[Measurability from maximal subset products]
\label{lem:measurable-ordered-tuple}
\leanok
\lean{ArakelovGeometry.measurable_perm_eval_of_measurable_iSup_prod}
\uses{lem:ordered-subset-product}
\source{chen-moriwaki-adelic:page-0198}
Let $\Omega$ be a measurable space and let
$a_0,\ldots,a_{n-1}:\Omega\to\mathbb R_{>0}$ be arbitrary functions.
For each $\omega$, choose a permutation $\pi_\omega$ making
$a_{\pi_\omega(0)}(\omega)\geq\cdots\geq a_{\pi_\omega(n-1)}(\omega)$.
Suppose that for every $0\leq k\leq n$ the function
\[
 P_k(\omega)=\max_{|J|=k}\prod_{j\in J}a_j(\omega)
\]
is measurable. Then every function
$\omega\mapsto a_{\pi_\omega(i)}(\omega)$ is measurable.
\end{lemma}
\begin{proof}
Lemma~\ref{lem:ordered-subset-product} identifies $P_k$ with the product
of the first $k$ entries in the chosen ordering. These products are
strictly positive, including $P_0=1$. Hence
$a_{\pi_\omega(i)}(\omega)=P_{i+1}(\omega)/P_i(\omega)$.
A quotient of measurable real functions is measurable. No measurability
of the permutations is used.
\end{proof}

\begin{lemma}[Measurable maximal products of embedding values]
\label{lem:measurable-embedding-products}
\leanok
\lean{ArakelovGeometry.AdelicStructure.measurable_iSup_prod_embeddings}
\uses{lem:orbit-maximum}
\source{chen-moriwaki-adelic:page-0198}
Let $E/K$ be finite separable, let $b\in E$, and choose for each
$\omega$ an extension $v_\omega$ of $\phi_\omega$ to an algebraic
closure $C$ of $K$. Write $S=\operatorname{Hom}_K(E,C)$.
For every $0\leq k\leq\#S$, the function
\[
 \omega\longmapsto \max_{J\subseteq S,\,\#J=k}
       \prod_{\sigma\in J}v_\omega(\sigma(b))
\]
is measurable.
\end{lemma}
\begin{proof}
For such a subset $J$ put $c_J=\prod_{\sigma\in J}\sigma(b)$.
Each $\sigma(b)$ is separable, so $c_J$ is separable because the
separable elements in $C$ form a field. Every $\tau\in\operatorname{Aut}_K(C)$
permutes $S$ by postcomposition, with inverse given by postcomposition
with $\tau^{-1}$. Thus $\tau(c_J)=c_{\tau J}$, where $\#\tau J=k$.
It follows that
\[
 \max_{\#J=k}v_\omega(c_J)
 =\max_{\#J=k}\sup_\tau v_\omega(\tau(c_J)).
\]
Indeed each term on the right is bounded by the left, and the identity
automorphism gives the reverse inequality. There are finitely many $J$,
and each inner supremum is measurable by Lemma~\ref{lem:orbit-maximum}.
Their maximum is therefore measurable. Multiplicativity of $v_\omega$
identifies this maximum with the stated function.
\end{proof}

\begin{lemma}[Measurable ordered conjugate values]\label{lem:ordered-conjugate-values}
\leanok
\source{chen-moriwaki-adelic:page-0198}
\uses{lem:measurable-embedding-products,lem:measurable-ordered-tuple}
\lean{ArakelovGeometry.AdelicStructure.measurable_ordered_abv_embeddings}
Let $E/K$ be finite separable and $b\in E\setminus\{0\}$. Enumerate
the embeddings $\sigma:E\hookrightarrow C$ and, for each $\omega$,
permute their values $v_\omega(\sigma(b))$ into decreasing order
$s_1(\omega)\geq\cdots\geq s_m(\omega)>0$, where $v_\omega$ is any
extension of $\phi_\omega$ to an algebraic closure $C$. Each $s_i$ is
measurable. For $E=K(b)$ these are the ordered absolute values of the
distinct conjugates of $b$; the general statement allows repetitions.
\end{lemma}
\begin{proof}
Every embedding is injective, so every $\sigma(b)$ is nonzero and its
absolute value is positive. Lemma~\ref{lem:measurable-embedding-products}
gives measurability of all maximal subset products. After transporting
subsets along the chosen enumeration, Lemma~\ref{lem:measurable-ordered-tuple}
gives the assertion. This is the ordered-product argument of the reference,
applied to the entire embedding list rather than only to a simple generator.
\end{proof}

\begin{lemma}[Measurable sums of embedding values]
\label{lem:measurable-embedding-sum}
\leanok
\lean{ArakelovGeometry.AdelicStructure.measurable_sum_abv_embeddings}
\uses{lem:ordered-conjugate-values}
\source{chen-moriwaki-adelic:page-0198,chen-moriwaki-adelic:page-0199}
Under the preceding choices, for every $b\in E$ the function
$\omega\mapsto\sum_{\sigma:E\hookrightarrow C}v_\omega(\sigma(b))$
is measurable.
\end{lemma}
\begin{proof}
For $b=0$ the sum is zero. For $b\ne0$, choose a decreasing ordering
of the finite list of embedding values at each $\omega$.
Each ordered entry is measurable by Lemma~\ref{lem:ordered-conjugate-values}.
Their finite sum is measurable and equals the original sum, because a
permutation preserves a finite sum.
\end{proof}

\begin{lemma}[Evaluation integrals as embedding averages]
\label{lem:evaluation-embedding-average}
\leanok
\lean{ArakelovGeometry.AdelicStructure.ParameterSpaceExtension.fibreIntegral_eval_eq_sum_embeddings}
\uses{lem:finite-fibre-sums,lem:weighted-embedding-sums}
\source{chen-moriwaki-adelic:page-0198,chen-moriwaki-adelic:page-0199}
Fix $\omega$ and any extension $v$ of $\phi_\omega$ to an algebraic
closure $C$ of $K$. For every $b\in L$,
\[
 I_{L/K}e_b(\omega)=\frac1{[L:K]}
   \sum_{\sigma:L\hookrightarrow C}v(\sigma(b)).
\]
\end{lemma}
\begin{proof}
The finite-sum formula gives
$I_{L/K}e_b(\omega)=[L:K]^{-1}\sum_{x\in X_\omega}
[L_x:K_\omega]|b|_x$. The points of $X_\omega$ correspond bijectively
to the absolute values on $L$ restricting to $\phi_\omega$: a point
has this absolute value and base coordinate $\omega$, which also
reconstruct it. Apply Lemma~\ref{lem:weighted-embedding-sums} to the
real test function $w\mapsto w(b)$. Its value at the pullback along
$\sigma$ is $v(\sigma(b))$, yielding the formula.
\end{proof}

\begin{lemma}[Measurable averages of evaluations]\label{lem:measurable-evaluation-average}
\leanok
\source{chen-moriwaki-adelic:page-0198,chen-moriwaki-adelic:page-0199}
\uses{lem:algebraic-absolute-value-extension,lem:evaluation-embedding-average,lem:measurable-embedding-sum}
\lean{ArakelovGeometry.AdelicStructure.ParameterSpaceExtension.measurable_fibreIntegral_eval}
For every $b\in L$, the function $I_{L/K}e_b$ is measurable.
\end{lemma}
\begin{proof}
Extend each base absolute value to an algebraic closure by
Lemma~\ref{lem:algebraic-absolute-value-extension}; no measurability
of these choices is needed. Lemma~\ref{lem:evaluation-embedding-average}
expresses $I_{L/K}e_b$ as $[L:K]^{-1}$ times the embedding-value sum.
That sum is measurable by Lemma~\ref{lem:measurable-embedding-sum},
so multiplication by this fixed real constant proves the claim.
\end{proof}

\subsection{The measurable-function extension argument}

\begin{lemma}[Evaluation products generate the extension sigma algebra]
\label{lem:evaluation-products-generate}
\leanok
\lean{ArakelovGeometry.AdelicStructure.ParameterSpaceExtension.measurableSpace_eq_iSup_comap_eval_mul_base}
\uses{def:extension-parameters,lem:extension-coordinates}
\source{chen-moriwaki-adelic:page-0199,chen-moriwaki-adelic:page-0200}
The sigma algebra $\mathcal A_L$ is generated by the functions
$e_b(\varphi\circ\pi)$, where $b\in L$ and
$\varphi:\Omega\to\mathbb R$ is measurable, with no boundedness
condition on the product or on $\varphi$.
\end{lemma}
\begin{proof}
Each product is measurable for $\mathcal A_L$, so the sigma algebra
they generate is contained in $\mathcal A_L$. Conversely take
$\varphi=1$ to recover every evaluation $e_b$. For any
$A\in\mathcal A$, take $b=1$ and $\varphi=\mathbf1_A$; the product
is $\mathbf1_{\pi^{-1}A}$. Its inverse image of $\{1\}$ is
$\pi^{-1}A$, so the projection is measurable for the generated sigma
algebra. The evaluations and the projection generate $\mathcal A_L$,
which proves the reverse inclusion. This is the generator argument
underlying the bounded family in the reference before bounded
localization is imposed.
\end{proof}

\begin{lemma}[A coefficient bound for roots]\label{lem:algebraic-root-bound}
\source{chen-moriwaki-adelic:page-0200}
\notready
If $b$ satisfies $b^m+\mu_1b^{m-1}+\cdots+\mu_m=0$ with $m\geq1$
and $\mu_i\in K$, then for $x\in X_\omega$,
\[
 |b|_x\leq m\max(1,|\mu_1|_\omega,\ldots,|\mu_m|_\omega).
\]
\end{lemma}
\begin{proof}
Put $r=|b|_x$ and $M=\max(1,|\mu_1|_\omega,\ldots,|\mu_m|_\omega)$.
If $r>mM$, then $r>1$. Dividing the triangle inequality for the polynomial
relation by $r^m$ gives
$1\leq\sum_{i=1}^m|\mu_i|_\omega/r^i\leq mM/r<1$, a contradiction.
\end{proof}

\begin{lemma}[Bounded multiplicative generators]\label{lem:bounded-evaluation-generators}
\source{chen-moriwaki-adelic:page-0199,chen-moriwaki-adelic:page-0200}
\uses{def:extension-parameters,def:adelic-structure,lem:algebraic-root-bound}
\notready
Let $C$ consist of the bounded nonnegative functions on $\Omega_L$ of
the form $e_b(\varphi\circ\pi)$, where $b\in L$ and $\varphi$ is a
bounded nonnegative measurable function on $\Omega$. Then $1\in C$,
$C$ is closed under multiplication, and $\sigma(C)=\mathcal A_L$.
\end{lemma}
\begin{proof}
Take $b=1$, $\varphi=1$ to obtain the constant one. Multiplicativity
of absolute values gives
$e_b(\varphi\circ\pi)e_c(\psi\circ\pi)=e_{bc}((\varphi\psi)\circ\pi)$;
the product is bounded and nonnegative. All these functions are
$\mathcal A_L$-measurable, so $\sigma(C)\subseteq\mathcal A_L$.
For every $A\in\mathcal A$, choosing $b=1$ and $\varphi=\mathbf1_A$
shows that $\mathbf1_{\pi^{-1}A}\in C$; hence the projection is
$\sigma(C)$-measurable.

Fix $b\in L$ and its monic minimal polynomial of degree $m$. Set
$A_N=\{\omega:\max_i|\mu_i|_\omega\leq N\}$ for positive integers $N$.
These sets are measurable and increase to $\Omega$. The preceding root
bound gives $0\leq e_b(\mathbf1_{A_N}\circ\pi)\leq m\max(1,N)$.
Thus each such function belongs to $C$, and their pointwise limit is
$e_b$. Every evaluation is therefore $\sigma(C)$-measurable. Together
with the projection these generate $\mathcal A_L$, proving equality.
\end{proof}

\begin{lemma}[Bounded functional monotone closure]\label{lem:functional-monotone-class}
\source{chen-moriwaki-adelic:page-0199}
\notready
Let $C$ be a multiplicatively closed family of bounded real functions
on a set $X$ containing $1$. Suppose a real vector space $H$ of bounded
functions contains $C$ and is closed under pointwise limits of uniformly
bounded sequences. Then $H$ contains every bounded $\sigma(C)$-measurable
real function.
\end{lemma}
\begin{proof}
This vector-space formulation supplies the functional monotone-class
step used in the proof of Theorem 3.3.4.
Let $D$ be the intersection of all vector spaces of bounded functions
containing $C$ and closed under such limits. Then $D\subseteq H$ and
$D$ contains the algebra $A$ spanned by $C$. For fixed $g\in A$, the
set $\{f\in D:fg\in D\}$ is a vector space with the same limit closure
and contains $A$. Its closure property follows because $g$ is bounded.
Minimality gives $fg\in D$ for all $f\in D$. Fixing now $f\in D$ and
applying the same argument to $\{g\in D:fg\in D\}$ shows that $D$ is
an algebra.

To see that $|f|\in D$, first rescale to $|f|\leq1$. Define
$q_0=0$ and $q_{j+1}=q_j+(f^2-q_j^2)/2$. These are polynomials in $f$.
Pointwise, $0\leq q_j\leq|f|$ and $q_j$ increases: if $0\leq q\leq t\leq1$,
then $0\leq(t^2-q^2)/2\leq t-q$. Its limit $q$ satisfies
$q=q+(f^2-q^2)/2$, hence $q=|f|$. The limit closure proves the claim;
rescaling removes the bound one. Thus positive parts and finite minima
belong to $D$. For $f\in D$ and $r\in\mathbb R$, the bounded functions
$\min(1,j(f-r)^+)$ converge to $\mathbf1_{\{f>r\}}$, which belongs to $D$.

Let $\mathcal E=\{B\subseteq X:\mathbf1_B\in D\}$. It contains $X$,
is closed under complements, and under finite intersections because
$D$ is an algebra. For a countable union, the indicators of its finite
unions belong to $D$, are bounded by one, and converge to the indicator
of the union. Hence $\mathcal E$ is a sigma algebra. It contains the
inverse images of rays under all members of $C$, so contains $\sigma(C)$.
Finally, a bounded measurable real function is the pointwise limit of
uniformly bounded finite-valued measurable functions obtained by rounding
down to multiples of $2^{-j}$. Their level-set indicators belong to $D$,
so the simple functions and their limit belong to $D\subseteq H$.
\end{proof}

\begin{lemma}[Functional closure under pointwise sequences]
\label{lem:sequential-functional-monotone-class}
\leanok
\source{chen-moriwaki-adelic:page-0199,chen-moriwaki-adelic:page-0200}
\lean{MeasureTheory.measurable_mem_of_isSeqClosed}
Let $C$ be a multiplicatively closed family of real functions on a set
$X$ containing $1$. Let $H$ be a real vector subspace of $\mathbb R^X$
containing $C$. Assume that whenever $f_n\in H$ and $f_n$ converges
pointwise to a real function $f$, one has $f\in H$. Then $H$ contains
every $\sigma(C)$-measurable real function.
\end{lemma}
\begin{proof}
Let $D$ be the intersection of all real vector subspaces containing
$C$ and closed under pointwise sequential limits. Then $D\subseteq H$,
and $D$ has the same limit closure. For $g\in C$, the vector space
$\{f\in D:fg\in D\}$ contains $C$, by multiplicative closure of $C$,
and is closed under pointwise sequential limits. Minimality implies
$fg\in D$ for all $f\in D$. Fix such an $f$. The vector space
$\{g\in D:fg\in D\}$ now contains $C$ and has the same limit closure.
Minimality again proves that $D$ is an algebra containing $1$.

For every continuous $u:\mathbb R\to\mathbb R$ and $f\in D$, choose,
by polynomial approximation on compact intervals, a polynomial $P_n$
with $|P_n(t)-u(t)|<1/(n+1)$ for $|t|\leq n$. For each $x$, this
estimate applies to $t=f(x)$ for all sufficiently large $n$, so
$P_n(f(x))\to u(f(x))$. Polynomial expressions in $f$ belong to $D$;
therefore $u\circ f\in D$.

For $r\in\mathbb R$, the continuous functions
$u_n(t)=\min(1,n\max(0,t-r))$ converge pointwise to
$\mathbf1_{(r,\infty)}(t)$. Hence
$\mathbf1_{\{f>r\}}\in D$ for every $f\in D$.
Let $\mathcal E=\{A\subseteq X:\mathbf1_A\in D\}$.
It contains $X$ and is closed under complements and finite
intersections because $D$ is an algebra containing $1$.
The indicators of the finite unions of a sequence of members of
$\mathcal E$ converge pointwise to the indicator of their union.
Thus $\mathcal E$ is a sigma algebra. It contains the inverse images
of open rays under every member of $C$, and consequently contains
$\sigma(C)$.
Every measurable real simple function is a finite linear combination
of such indicators and belongs to $D$. Finally, every real measurable
function is a pointwise limit of measurable simple functions, obtained
by truncating and rounding to a successively finer rational grid.
It belongs to $D\subseteq H$ by sequential closure.
\end{proof}

\begin{theorem}[Measurability of fibre integration]\label{thm:measurable-fibre-integral}
\leanok
\source{chen-moriwaki-adelic:page-0196,chen-moriwaki-adelic:page-0199,chen-moriwaki-adelic:page-0200}
\uses{lem:finite-fibre-limits,lem:fibre-additivity,lem:fibre-homogeneity,prop:fibre-base-multiplication,lem:measurable-evaluation-average,lem:evaluation-products-generate,lem:sequential-functional-monotone-class}
\lean{ArakelovGeometry.AdelicStructure.ParameterSpaceExtension.measurable_fibreIntegral}
If $f:\Omega_L\to\mathbb R$ is measurable, then $If:\Omega\to\mathbb R$
is measurable.
\end{theorem}
\begin{proof}
Let $H$ consist of all real functions whose fibre integrals are measurable.
Linearity makes $H$ a real vector space. If a sequence in $H$ converges
pointwise to a real function, Lemma~\ref{lem:finite-fibre-limits} gives
pointwise convergence of its fibre integrals to the fibre integral of
the limit. A pointwise sequential limit of measurable real functions
is measurable, so $H$ has the same closure.
Let $C$ be the family $e_b(\varphi\circ\pi)$ with $b\in L$ and
$\varphi:\Omega\to\mathbb R$ measurable. It contains $1$ by taking
$b=1$, $\varphi=1$. Multiplicativity of absolute values shows that
the product of the members associated to $(b,\varphi)$ and
$(c,\psi)$ is associated to $(bc,\varphi\psi)$, so $C$ is
multiplicatively closed. Multiplication by a base function gives
$I(e_b(\varphi\circ\pi))=\varphi Ie_b$, measurable by
Lemma~\ref{lem:measurable-evaluation-average}. Thus $C\subseteq H$.
Lemma~\ref{lem:evaluation-products-generate} identifies $\sigma(C)$
with $\mathcal A_L$. Lemma~\ref{lem:sequential-functional-monotone-class}
therefore places every measurable real function in $H$. The finiteness
of each fibre supplies the unrestricted sequential closure used here;
the bounded localization in the reference is unnecessary for this
specialization of its functional monotone-class argument.
\end{proof}

\begin{lemma}[Measurability of nonnegative fibre integration]\label{lem:measurable-fibre-lintegral}
\source{chen-moriwaki-adelic:page-0200,chen-moriwaki-adelic:page-0201}
\uses{thm:measurable-fibre-integral,lem:finite-fibre-integrability}
\lean{ArakelovGeometry.AdelicStructure.ParameterSpaceExtension.measurable_fibreLIntegral}
For every measurable $f:\Omega_L\to[0,\infty]$, the function $I^+f$
is measurable.
\end{lemma}
\begin{proof}
For a measurable set $A$, its real indicator has a measurable fibre
integral. This integral is $P_\omega(A\cap X_\omega)$, a finite
nonnegative number since $P_\omega$ is a probability measure.
Thus $\omega\mapsto P_\omega(A\cap X_\omega)$ is measurable also
as an extended nonnegative function. For any $c\in[0,\infty]$,
$I^+(c\mathbf1_A)=cP_\omega(A\cap X_\omega)$ is measurable.
Additivity of nonnegative integrals proves the assertion for finite
nonnegative measurable simple functions. Choose such functions
$s_j\uparrow f$. Monotone convergence on each fibre gives
$I^+f=\sup_j I^+s_j$, a measurable function.
\end{proof}

\subsection{The induced measure and integral identities}

\begin{definition}[The induced measure]\label{def:separable-induced-measure}
\source{chen-moriwaki-adelic:page-0200,chen-moriwaki-adelic:page-0201}
\uses{lem:measurable-fibre-lintegral,lem:finite-fibre-sums}
\lean{ArakelovGeometry.AdelicStructure.ParameterSpaceExtension.inducedMeasure,ArakelovGeometry.AdelicStructure.ParameterSpaceExtension.inducedMeasure_apply}
The induced measure $\nu_L$ on $(\Omega_L,\mathcal A_L)$ is characterized by
\[
 \nu_L(A)=\int_\Omega I^+\mathbf1_A\,d\nu
 \qquad(A\in\mathcal A_L).
\]
\end{definition}
\begin{proof}
The formula assigns zero to the empty set. For disjoint measurable
$A_j$, the finite-fibre formula gives
$I^+\mathbf1_{\bigcup_j A_j}=\sum_j I^+\mathbf1_{A_j}$.
Each summand is measurable by Lemma~\ref{lem:measurable-fibre-lintegral}.
Monotone convergence applied to partial sums gives countable additivity
of the displayed set function. It therefore defines a measure on
$\mathcal A_L$, with the stated values.
\end{proof}

\begin{theorem}[Nonnegative integration formula]\label{thm:separable-lintegral}
\source{chen-moriwaki-adelic:page-0201}
\uses{def:separable-induced-measure,lem:measurable-fibre-lintegral,lem:finite-fibre-sums}
\lean{ArakelovGeometry.AdelicStructure.ParameterSpaceExtension.lintegral_inducedMeasure}
For every measurable $f:\Omega_L\to[0,\infty]$,
\[
 \int_{\Omega_L}f\,d\nu_L=\int_\Omega I^+f\,d\nu.
\]
\end{theorem}
\begin{proof}
For an indicator this is the definition of $\nu_L$. Finite additivity
and nonnegative homogeneity of both integrals prove the formula for
nonnegative measurable simple functions. Choose simple functions
$s_j\uparrow f$. Monotone convergence on $\Omega_L$ gives convergence
of the left integrals. Fibrewise monotone convergence gives
$I^+s_j\uparrow I^+f$, and these are measurable functions. Monotone
convergence on $\Omega$ gives convergence of the right integrals.
Taking the limit in the simple-function identities proves the assertion.
\end{proof}

\begin{theorem}[Integrability criterion]\label{thm:separable-integrability}
\source{chen-moriwaki-adelic:page-0201}
\uses{thm:separable-lintegral,thm:measurable-fibre-integral,lem:finite-fibre-sums}
\lean{ArakelovGeometry.AdelicStructure.ParameterSpaceExtension.integrable_inducedMeasure_iff}
A measurable real function $f$ is $\nu_L$-integrable if and only if
$I|f|$ is $\nu$-integrable.
\end{theorem}
\begin{proof}
Both $f$ and $I|f|$ are measurable. The latter is nonnegative and finite
at every base point by the finite-sum formula. The nonnegative integration
formula gives
\[
 \int_{\Omega_L}|f|\,d\nu_L=\int_\Omega I|f|\,d\nu
\]
as an equality in $[0,\infty]$. Integrability on either side is precisely
finiteness of this number.
\end{proof}

\begin{lemma}[Null sets and finite fibres]\label{lem:finite-fibre-ae}
\leanok
\uses{lem:separable-ae-fibres,lem:finite-fibre-sum}
\source{chen-moriwaki-adelic:page-0201,chen-moriwaki-adelic:page-0202}
\lean{ArakelovGeometry.AdelicStructure.ParameterSpaceExtension.ae_fibreProbabilityMeasure_iff,ArakelovGeometry.AdelicStructure.ParameterSpaceExtension.ae_forall_fibre_of_ae,ArakelovGeometry.AdelicStructure.ParameterSpaceExtension.fibreIntegral_congr_ae}
On a finite fibre, a property holds $P_\omega$-almost everywhere if
and only if it holds at every point. Consequently a property holding
$\nu_L$-almost everywhere holds at every point of $X_\omega$ for
$\nu$-almost every $\omega$.
If two real functions $f,g$ on $\Omega_L$ agree $\nu_L$-almost everywhere,
then for $\nu$-almost every $\omega$ they agree at every point of
$X_\omega$. In particular $If=Ig$ almost everywhere.
\end{lemma}
\begin{proof}
Each singleton $\{x\}$ has mass $p_\omega(x)>0$. A null exceptional
set therefore contains no point, proving the first assertion; the
converse is immediate. Lemma~\ref{lem:separable-ae-fibres} gives a
fibre-almost-everywhere property on almost every fibre, and the first
assertion turns it into an everywhere property on those fibres.
Apply this to $f=g$. The integrands agree at every point of almost
every fibre, so their integrals agree on those same fibres.
\end{proof}

\begin{lemma}[Almost everywhere statements on fibres]\label{lem:separable-ae-fibres}
\source{chen-moriwaki-adelic:page-0201,chen-moriwaki-adelic:page-0202}
\uses{def:separable-induced-measure,lem:measurable-fibre-lintegral}
\lean{ArakelovGeometry.AdelicStructure.ParameterSpaceExtension.ae_ae_fibre_of_ae}
If a property holds for $\nu_L$-almost every $x\in\Omega_L$, then
for $\nu$-almost every $\omega$ it holds for $P_\omega$-almost every
$x\in X_\omega$.
\end{lemma}
\begin{proof}
The exceptional set is contained in a measurable $\nu_L$-null set $N$.
The defining formula for $\nu_L$ gives
$\int_\Omega P_\omega(N\cap X_\omega)\,d\nu=0$.
The integrand is measurable and nonnegative, so it is zero almost
everywhere. Outside this exceptional base set, the original property
fails only on the $P_\omega$-null set $N\cap X_\omega$.
\end{proof}

\begin{theorem}[Real integration formula]\label{thm:separable-integral}
\source{chen-moriwaki-adelic:page-0201,chen-moriwaki-adelic:page-0202}
\uses{thm:separable-lintegral,thm:separable-integrability,lem:separable-ae-fibres,lem:fibre-additivity,lem:fibre-homogeneity,thm:measurable-fibre-integral}
\lean{ArakelovGeometry.AdelicStructure.ParameterSpaceExtension.integral_inducedMeasure}
For every $\nu_L$-integrable real function $f$,
\[
 \int_{\Omega_L}f\,d\nu_L=\int_\Omega If\,d\nu.
\]
The integrability assumption allows almost everywhere measurable functions.
\end{theorem}
\begin{proof}
First suppose $f$ is measurable. Its positive and negative parts $f^+,f^-$
have finite integrals. The nonnegative integration formula applies to both.
Moreover $If^+$ and $If^-$ are integrable, since they are measurable and
their nonnegative integrals are finite. Subtracting their integral
identities, and using $If=If^+-If^-$, gives the asserted formula; it
also proves integrability of $If$.

For an arbitrary integrable $f$, choose a measurable real representative
$g=f$ almost everywhere. Such a representative is obtained from the
measurable simple functions witnessing almost everywhere strong
measurability, taking their limit where it exists and zero elsewhere.
It is integrable with the same integral as $f$. By
Lemma~\ref{lem:separable-ae-fibres}, $If=Ig$ almost everywhere. Thus $If$
is integrable and has the same integral as $Ig$. Apply the measurable
case to $g$ to conclude. This extends the measurable-integrand formulation
of Theorem 3.3.7(2) by equality almost everywhere.
\end{proof}

\begin{proposition}[Pushforward to the base]\label{prop:separable-pushforward}
\source{chen-moriwaki-adelic:page-0201,chen-moriwaki-adelic:page-0202}
\uses{def:extension-parameters,def:separable-induced-measure,lem:finite-fibre-sums}
\lean{ArakelovGeometry.AdelicStructure.ParameterSpaceExtension.map_inducedMeasure_proj}
The projection satisfies $\pi_*\nu_L=\nu$.
\end{proposition}
\begin{proof}
For measurable $A\subseteq\Omega$, the indicator of $\pi^{-1}A$ is
constant on $X_\omega$ with value $\mathbf1_A(\omega)$. Since the fibre
measure has total mass one, $I^+\mathbf1_{\pi^{-1}A}=\mathbf1_A$.
The definition of $\nu_L$ yields $\nu_L(\pi^{-1}A)=\nu(A)$.
Equality on measurable sets is equality of measures.
\end{proof}

\begin{lemma}[Nonnegative integrals and integrability of base pullbacks]\label{lem:separable-pullback-integrability}
\leanok
\source{chen-moriwaki-adelic:page-0201,chen-moriwaki-adelic:page-0202}
\uses{prop:separable-pushforward,def:extension-parameters}
\lean{ArakelovGeometry.AdelicStructure.ParameterSpaceExtension.lintegral_comp_proj_inducedMeasure,ArakelovGeometry.AdelicStructure.ParameterSpaceExtension.integrable_comp_proj_inducedMeasure}
For every measurable $\varphi:\Omega\to[0,\infty]$,
\[
 \int_{\Omega_L}\varphi\circ\pi\,d\nu_L=\int_\Omega\varphi\,d\nu.
\]
For every $\nu$-integrable real function $\varphi$, its pullback
$\varphi\circ\pi$ is $\nu_L$-integrable.
\end{lemma}
\begin{proof}
For measurable nonnegative functions, the integral transformation formula
under a measurable map follows first for indicators, then for simple
functions, and then by monotone convergence. Apply it to $|\varphi|$
and use $\pi_*\nu_L=\nu$ to obtain integrability. If $\varphi$
is only almost everywhere measurable, replace it by a measurable
representative. The inverse image of a measurable base null set is null
by the pushforward identity, so integrability remains true for the original
function.
\end{proof}

\begin{corollary}[Integration of base pullbacks]\label{cor:separable-pullback-integral}
\source{chen-moriwaki-adelic:page-0201,chen-moriwaki-adelic:page-0202}
\uses{lem:separable-pullback-integrability,prop:separable-pushforward,def:extension-parameters}
\lean{ArakelovGeometry.AdelicStructure.ParameterSpaceExtension.integral_comp_proj_inducedMeasure}
For every $\nu$-integrable real function $\varphi$,
\[
 \int_{\Omega_L}\varphi\circ\pi\,d\nu_L=\int_\Omega\varphi\,d\nu.
\]
\end{corollary}
\begin{proof}
For a measurable $\varphi$, apply the preceding lemma to its positive
and negative parts and subtract their finite integral identities.
For an almost everywhere measurable $\varphi$, choose a measurable
representative; its pullback equals $\varphi\circ\pi$ almost everywhere
because the pushforward identity takes base null sets to null inverse
images. Replacing integrands by these equal representatives gives the
same equality.
\end{proof}

\subsection{Logarithms, norms, and properness}

\begin{lemma}[The fibrewise logarithmic norm identity]\label{lem:fibre-log-norm}
\leanok
\source{chen-moriwaki-adelic:page-0202}
\uses{prop:local-norm-product,lem:finite-fibre-sums}
\lean{ArakelovGeometry.AdelicStructure.ParameterSpaceExtension.fibreIntegral_log_eq_log_norm}
For $b\in L^\times$,
\[
 I_{L/K}(\log e_b)(\omega)
 =\frac1n\log|N_{L/K}(b)|_\omega.
\]
\end{lemma}
\begin{proof}
The local norm product formula gives
$|N_{L/K}(b)|_\omega=\prod_{x\in X_\omega}|b|_x^{[L_x:K_\omega]}$.
Every factor is positive. Taking logarithms and dividing by $n$
turns this product into the finite-sum formula for the fibre integral.
The operator is $I_{L/K}$, with domain functions on $\Omega_L$; this
fixes the reversed subscripts in the printed formula (3.10).
\end{proof}

\begin{lemma}[Root bounds from bounds on finite sums]\label{lem:root-bound-from-sums}
\leanok
\source{chen-moriwaki-adelic:page-0203,chen-moriwaki-adelic:page-0204}
\lean{AbsoluteValue.apply_root_le_of_sum_bound}
Let $P=T^m+\sum_{i<m}\mu_iT^i$ be monic over $K$ and let $a\in L$
be a root in an extension field, endowed with an absolute value $w$.
Suppose $A,M\geq1$, $w(\mu_i)\leq M$ for $i<m$, and every family
of $m$ elements with absolute values at most $B\geq0$ has sum of
absolute value at most $AB$. Then $w(a)\leq AM$.
\end{lemma}
\begin{proof}
The existence of a root forces $m>0$. Put $r=w(a)$ and suppose
$r>AM\geq1$. For $i<m$ one has $w(\mu_i a^i)\leq Mr^{m-1}$.
Applying the assumed finite-sum bound to
$a^m=-\sum_{i<m}\mu_i a^i$ gives $r^m\leq AMr^{m-1}$.
Since $r^{m-1}>0$, division gives $r\leq AM$, a contradiction.
\end{proof}

\begin{lemma}[Nonarchimedean root bound]\label{lem:nonarchimedean-root-bound}
\leanok
\uses{lem:root-bound-from-sums}
\source{chen-moriwaki-adelic:page-0202,chen-moriwaki-adelic:page-0203}
\lean{AbsoluteValue.apply_root_le_of_isNonarchimedean}
For a monic polynomial $P=T^m+\sum_{i<m}\mu_iT^i$ over $K$,
a root $a$ in an extension field carrying a nonarchimedean absolute
value satisfies $w(a)\leq M$ whenever $M\geq1$ and
$w(\mu_i)\leq M$ for every $i<m$.
\end{lemma}
\begin{proof}
The ultrametric inequality, applied inductively, bounds the absolute
value of a nonempty finite sum by the largest absolute value of a
summand. Apply Lemma~\ref{lem:root-bound-from-sums} with $A=1$.
\end{proof}

\begin{lemma}[Archimedean root bound]\label{lem:archimedean-root-bound}
\leanok
\uses{lem:root-bound-from-sums,lem:powered-complex-sum-bound}
\source{chen-moriwaki-adelic:page-0203,chen-moriwaki-adelic:page-0204}
\lean{AbsoluteValue.apply_root_le_of_complex_embedding}
Suppose $w=|\sigma(\cdot)|^s$ for a complex embedding and $s\geq0$.
For a root $a$ of a monic polynomial
$P=T^m+\sum_{i<m}\mu_iT^i$, if $M\geq1$ and $w(\mu_i)\leq M$
for $i<m$, then $w(a)\leq m^sM$.
\end{lemma}
\begin{proof}
The root forces $m\geq1$, hence $m^s\geq1$.
Lemma~\ref{lem:powered-complex-sum-bound} bounds each sum of $m$
terms of absolute value at most $B$ by $m^sB$.
Apply Lemma~\ref{lem:root-bound-from-sums} with $A=m^s$.
\end{proof}

\begin{definition}[Coefficient logarithmic maximum]\label{def:coefficient-log-envelope}
\leanok
\source{chen-moriwaki-adelic:page-0203}
\uses{def:adelic-structure}
\lean{ArakelovGeometry.AdelicStructure.coefficientLogBound}
For a polynomial $P\in K[T]$ of degree $m$, put
$G_P(\omega)=\max_{0\leq i\leq m}\ell_i(\omega)$, where
$\ell_i(\omega)=\log|P_i|_\omega$ if $P_i\ne0$ and
$\ell_i(\omega)=0$ if $P_i=0$.
\end{definition}

\begin{lemma}[Regularity of the coefficient maximum]\label{lem:coefficient-log-envelope}
\leanok
\source{chen-moriwaki-adelic:page-0203}
\uses{def:coefficient-log-envelope}
\lean{ArakelovGeometry.AdelicStructure.measurable_coefficientLogBound,ArakelovGeometry.AdelicStructure.integrable_coefficientLogBound,ArakelovGeometry.AdelicStructure.coefficient_log_le,ArakelovGeometry.AdelicStructure.coefficientLogBound_nonneg}
The function $G_P$ is measurable and integrable, and it bounds every
$\ell_i$ with $i\leq m$ from above. If $P$ is monic, then $G_P\geq0$
and it equals the maximum of $0$ and the logarithmic absolute values
of the nonzero lower coefficients.
\end{lemma}
\begin{proof}
Each $\ell_i$ is either zero or measurable and integrable by the adelic
hypothesis. A finite maximum of such functions is measurable and
integrable, since the maximum of two real integrable functions is
integrable and the index set is nonempty. The defining maximum bounds
every entry. For monic $P$, the leading coefficient is $1$ and its
logarithm is $0$, which proves the last assertions.
\end{proof}

\begin{lemma}[An integrable upper bound for logarithmic evaluations]\label{lem:log-root-bound}
\leanok
\lean{ArakelovGeometry.AdelicStructure.measurable_logarithmicRootBound,ArakelovGeometry.AdelicStructure.integrable_logarithmicRootBound,ArakelovGeometry.AdelicStructure.logarithmicRootBound_nonneg,ArakelovGeometry.AdelicStructure.log_apply_le_rootBound}
\source{chen-moriwaki-adelic:page-0202,chen-moriwaki-adelic:page-0203,chen-moriwaki-adelic:page-0204}
\uses{def:adelic-structure,lem:archimedean-bound,lem:archimedean-extension,lem:archimedean-root-bound,lem:nonarchimedean-root-bound,lem:coefficient-log-envelope,lem:nonarchimedean-restriction,lem:archimedean-complex-embedding}
Let $b\in L^\times$ have minimal polynomial
$T^m+\mu_1T^{m-1}+\cdots+\mu_m$. There is a nonnegative measurable
integrable function $h$ on $\Omega$ such that $\log|b|_x\leq h(\pi x)$.
More precisely, with $\kappa$ the Archimedean exponent, extended by zero
at non-Archimedean places, one may take
\[
 g(\omega)=\max\bigl(\{0\}\cup
       \{\log|\mu_i|_\omega:\mu_i\ne0\}\bigr),\qquad
 h=(\log m)\kappa+g.
\]
\end{lemma}
\begin{proof}
The finite maximum $g$ is measurable and integrable: it is nonnegative
and at most $\sum_{\mu_i\ne0}|\log|\mu_i|_\omega|$, an integrable
function by the adelic hypothesis. The function $\kappa$ is measurable,
integrable, and nonnegative, so the same is true of $h$.

Put $M=\max(1,|\mu_1|_\omega,\ldots,|\mu_m|_\omega)$, so $\log M=g$.
At a non-Archimedean place, if $r=|b|_x>M$, the polynomial relation gives
$1\leq\max_i|\mu_i|_\omega/r^i<1$, a contradiction. Thus $r\leq M$.
At an Archimedean place put $k=\kappa(\omega)>0$. The normalized
absolute value $|\cdot|_x^{1/k}$ satisfies the triangle inequality.
If $r>m^kM$, then $r^{1/k}>mM^{1/k}\geq1$ and the polynomial relation
gives
\[
 1\leq\sum_{i=1}^m\frac{|\mu_i|_\omega^{1/k}}{r^{i/k}}
 \leq\frac{mM^{1/k}}{r^{1/k}}<1,
\]
again a contradiction. Hence in both cases $\log r\leq (\log m)k+g$,
where $k=0$ in the non-Archimedean case. Zero coefficients are omitted
from logarithms; they contribute nothing to $M$ because $1$ is included.
\end{proof}

\begin{theorem}[Integrability of logarithmic evaluations]\label{thm:separable-log-integrable}
\leanok
\source{chen-moriwaki-adelic:page-0202,chen-moriwaki-adelic:page-0203,chen-moriwaki-adelic:page-0204}
\uses{def:extension-parameters,def:adelic-structure,lem:log-root-bound,lem:fibre-log-norm,thm:separable-integrability,lem:separable-pullback-integrability,cor:fibre-pullback,lem:fibre-additivity,lem:fibre-homogeneity}
\lean{ArakelovGeometry.AdelicStructure.ParameterSpaceExtension.integrable_log_inducedMeasure}
For every $b\in L^\times$, $x\mapsto\log|b|_x$ is $\nu_L$-integrable.
\end{theorem}
\begin{proof}
Write $\ell=\log e_b$ and choose $h$ from Lemma~\ref{lem:log-root-bound}.
The function $u=h\circ\pi-\ell$ is nonnegative and measurable. Its fibre
integral is
\[
 Iu=h-I\ell=h-\frac1n\log|N_{L/K}(b)|.
\]
The norm of a nonzero element is nonzero, so the last logarithm is
integrable on the base. Thus $Iu=I|u|$ is integrable, and the integrability
criterion proves $u$ integrable. The pullback $h\circ\pi$ is integrable
by Lemma~\ref{lem:separable-pullback-integrability}. Their difference
$\ell=h\circ\pi-u$ is therefore integrable.
\end{proof}

\begin{definition}[The finite separable extension adelic structure]\label{def:separable-extension}
\leanok
\source{chen-moriwaki-adelic:page-0201,chen-moriwaki-adelic:page-0202,chen-moriwaki-adelic:page-0204}
\uses{def:adelic-structure,def:extension-parameters,def:separable-induced-measure,thm:separable-log-integrable}
\lean{ArakelovGeometry.AdelicStructure.finiteSeparableExtension}
The finite separable extension of $S$ to $L$ is the adelic structure
$S_L=(\Omega_L,\mathcal A_L,\nu_L,x\mapsto|\cdot|_x)$.
\end{definition}
\begin{proof}
The induced set function is a measure. Every nonzero logarithmic evaluation
is measurable by the definition of $\mathcal A_L$ and integrable by
Theorem~\ref{thm:separable-log-integrable}. These are exactly the adelic
structure conditions.
\end{proof}

\begin{proposition}[The global logarithmic norm identity]\label{prop:separable-norm-integral}
\leanok
\source{chen-moriwaki-adelic:page-0201,chen-moriwaki-adelic:page-0202,chen-moriwaki-adelic:page-0204}
\uses{thm:separable-log-integrable,thm:separable-integral,lem:fibre-log-norm,def:adelic-structure}
\lean{ArakelovGeometry.AdelicStructure.ParameterSpaceExtension.finrank_mul_integral_log_inducedMeasure}
For every $b\in L^\times$,
\[
 n\int_{\Omega_L}\log|b|_x\,d\nu_L(x)
 =\int_\Omega\log|N_{L/K}(b)|_\omega\,d\nu(\omega).
\]
\end{proposition}
\begin{proof}
Apply the real integration formula to the integrable function $\log e_b$.
Substitute the fibrewise logarithmic norm identity on its right side.
The base logarithm is integrable, so the factor $1/n$ passes through
the integral. Multiply by the positive integer $n$.
\end{proof}

\begin{proposition}[Preservation of properness]\label{prop:separable-proper}
\leanok
\source{chen-moriwaki-adelic:page-0201,chen-moriwaki-adelic:page-0204}
\uses{def:separable-extension,prop:separable-norm-integral,def:adelic-structure}
\lean{ArakelovGeometry.AdelicStructure.ParameterSpaceExtension.isProper_finiteSeparableExtension}
If $S$ is proper, then its finite separable extension $S_L$ is proper.
\end{proposition}
\begin{proof}
For $b\in L^\times$, its norm is in $K^\times$. Properness of $S$
sets the right side of the global logarithmic norm identity to zero.
Since $[L:K]>0$, the integral of $\log|b|$ on $\Omega_L$ is zero.
This is the product formula for every nonzero element of $L$, which
is precisely properness of $S_L$.
\end{proof}
